\documentclass[11pt]{article}
\usepackage[margin=1.1in]{geometry}
\usepackage{amsmath,amssymb,amsthm,mathtools,booktabs,array,tabularx,longtable}
\usepackage[dvipsnames]{xcolor}
\usepackage[colorlinks=true,linkcolor=MidnightBlue,citecolor=MidnightBlue,urlcolor=MidnightBlue]{hyperref}
\usepackage{fancyhdr}
\pagestyle{fancy}\fancyhf{}\renewcommand{\headrulewidth}{0pt}
\fancyfoot[R]{\thepage}
\fancyfoot[L]{\scriptsize\color{gray}Generated by the Transport Lab from approved parts of its store; not a deposited paper. Papers of record: \href{https://zenodo.org/communities/transport-papers}{Zenodo, Transport Papers}}
\fancypagestyle{plain}{\fancyhf{}\fancyfoot[R]{\thepage}\fancyfoot[L]{\scriptsize\color{gray}Generated by the Transport Lab from approved parts of its store; not a deposited paper. Papers of record: \href{https://zenodo.org/communities/transport-papers}{Zenodo, Transport Papers}}\renewcommand{\headrulewidth}{0pt}}
\newtheorem{theorem}{Theorem}[section]
\newtheorem{proposition}[theorem]{Proposition}
\newtheorem{lemma}[theorem]{Lemma}
\newtheorem{corollary}[theorem]{Corollary}
\newtheorem{conjecture}[theorem]{Conjecture}
\newtheorem{refutation}[theorem]{Refutation}
\newtheorem{prediction}[theorem]{Prediction}
\newtheorem{principle}[theorem]{Principle}
\newtheorem{observation}[theorem]{Observation}
\newtheorem{axiom}[theorem]{Axiom}
\theoremstyle{definition}
\newtheorem{definition}[theorem]{Definition}
\newtheorem{fact}[theorem]{Fact}
\newtheorem{assembled}[theorem]{Assembly}
\newtheorem{postulate}[theorem]{Postulate}
\newtheorem{premise}[theorem]{Premise}
\newtheorem{hypothesis}[theorem]{Hypothesis}
\newtheorem{remark}[theorem]{Remark}
\newtheorem{example}[theorem]{Example}
\theoremstyle{plain}
\setlength{\marginparwidth}{0.8in}
\newcommand{\changedmark}[2]{\marginpar{\raggedright\scriptsize\color{#1}#2}}
\providecommand{\Aa}{}\renewcommand{\Aa}{\mathcal{A}}
\providecommand{\Ad}{}\renewcommand{\Ad}{\widetilde{\mathrm{Ad}}}
\providecommand{\Aev}{}\renewcommand{\Aev}{\mathcal{A}^{+}}
\providecommand{\Aod}{}\renewcommand{\Aod}{\mathcal{A}^{-}}
\providecommand{\B}{}\renewcommand{\B}{\mathrm{B}}
\providecommand{\C}{}\renewcommand{\C}{\mathbb{C}}
\providecommand{\CAR}{}\renewcommand{\CAR}{\mathrm{CAR}}
\providecommand{\CS}{}\renewcommand{\CS}{{\mathbb{S}^{+}}}
\providecommand{\CSm}{}\renewcommand{\CSm}{{\mathbb{S}^{-}}}
\providecommand{\Cl}{}\renewcommand{\Cl}{\mathrm{Cl}}
\providecommand{\Con}{}\renewcommand{\Con}{\mathcal{C}}
\providecommand{\Dl}{}\renewcommand{\Dl}{\Delta}
\providecommand{\EN}{}\renewcommand{\EN}{E_{N}}
\providecommand{\End}{}\renewcommand{\End}{\mathrm{End}}
\providecommand{\F}{}\renewcommand{\F}{\mathbb{F}}
\providecommand{\G}{}\renewcommand{\G}{\Gamma}
\providecommand{\HH}{}\renewcommand{\HH}{\mathbb{H}}
\providecommand{\IS}{}\renewcommand{\IS}{I\sigma_3}
\providecommand{\J}{}\renewcommand{\J}{\mathbf{J}}
\providecommand{\LH}{}\renewcommand{\LH}{L_{H}}
\providecommand{\Mat}{}\renewcommand{\Mat}{\mathrm{M}}
\providecommand{\Mo}{}\renewcommand{\Mo}{M^{\circ}}
\providecommand{\Mp}{}\renewcommand{\Mp}{M^{+}}
\providecommand{\Msun}{}\renewcommand{\Msun}{M_{\odot}}
\providecommand{\Mt}{}\renewcommand{\Mt}{\widetilde{M}}
\providecommand{\Ncone}{}\renewcommand{\Ncone}{\mathcal{N}}
\providecommand{\Neff}{}\renewcommand{\Neff}{N_{\mathrm{eff}}}
\providecommand{\PCS}{}\renewcommand{\PCS}{\mathbb{P}{\mathbb{S}^{+}}}
\providecommand{\PT}{}\renewcommand{\PT}{\mathbb{PT}}
\providecommand{\Pin}{}\renewcommand{\Pin}{\mathrm{Pin}}
\providecommand{\Q}{}\renewcommand{\Q}{\mathbb{Q}}
\providecommand{\R}{}\renewcommand{\R}{\mathbb{R}}
\providecommand{\Ric}{}\renewcommand{\Ric}{\mathrm{Ric}}
\providecommand{\SBH}{}\renewcommand{\SBH}{S_{\mathrm{BH}}}
\providecommand{\Sc}{}\renewcommand{\Sc}{\mathbb{S}}
\providecommand{\Sch}{}\renewcommand{\Sch}{\mathrm{P}}
\providecommand{\Sig}{}\renewcommand{\Sig}{\Sigma}
\providecommand{\Som}{}\renewcommand{\Som}{S_\omega}
\providecommand{\Sp}{}\renewcommand{\Sp}{\mathbb{S}}
\providecommand{\Sph}{}\renewcommand{\Sph}{\mathrm{S}}
\providecommand{\Spi}{}\renewcommand{\Spi}{S_\pi}
\providecommand{\Spin}{}\renewcommand{\Spin}{\mathrm{Spin}}
\providecommand{\Sth}{}\renewcommand{\Sth}{S_{\mathrm{th}}}
\providecommand{\T}{}\renewcommand{\T}{\mathbb{T}}
\providecommand{\TH}{}\renewcommand{\TH}{T_{\mathrm{H}}}
\providecommand{\Tr}{}\renewcommand{\Tr}{\operatorname{Tr}}
\providecommand{\UHF}{}\renewcommand{\UHF}{\mathrm{UHF}}
\providecommand{\Z}{}\renewcommand{\Z}{\mathbb{Z}}
\providecommand{\Zc}{}\renewcommand{\Zc}{\mathrm{Z}}
\providecommand{\cmg}{}\renewcommand{\cmg}{\mathrm{cm^{2}\,g^{-1}}}
\providecommand{\dS}{}\renewcommand{\dS}{\mathrm{dS}}
\providecommand{\dsc}{}\renewcommand{\dsc}{\mathrm{dsc}}
\providecommand{\dtPG}{}\renewcommand{\dtPG}{dt}
\providecommand{\dtS}{}\renewcommand{\dtS}{dt_{\mathrm{S}}}
\providecommand{\ellv}{}\renewcommand{\ellv}{\ell}
\providecommand{\er}{}\renewcommand{\er}{e_r}
\providecommand{\et}{}\renewcommand{\et}{e_t}
\providecommand{\ev}{}\renewcommand{\ev}{\mathrm{ev}}
\providecommand{\gl}{}\renewcommand{\gl}{\mathfrak{gl}}
\providecommand{\hatV}{}\renewcommand{\hatV}{\widehat{V}}
\providecommand{\hb}{}\renewcommand{\hb}{\bar{h}}
\providecommand{\hf}{}\renewcommand{\hf}{\hat{f}}
\providecommand{\id}{}\renewcommand{\id}{\mathrm{id}}
\providecommand{\kpc}{}\renewcommand{\kpc}{\mathrm{kpc}}
\providecommand{\lP}{}\renewcommand{\lP}{\ell_{P}}
\providecommand{\mD}{}\renewcommand{\mD}{\mathbf{m}_{D}}
\providecommand{\mb}{}\renewcommand{\mb}{\bar m}
\providecommand{\mnu}{}\renewcommand{\mnu}{\Sigma m_{\nu}}
\providecommand{\nInf}{}\renewcommand{\nInf}{n_\infty}
\providecommand{\nO}{}\renewcommand{\nO}{n_0}
\providecommand{\ninf}{}\renewcommand{\ninf}{n_{\infty}}
\providecommand{\nz}{}\renewcommand{\nz}{n_{0}}
\providecommand{\proj}{}\renewcommand{\proj}{\operatorname{proj}}
\providecommand{\pv}{}\renewcommand{\pv}{\pi}
\providecommand{\rad}{}\renewcommand{\rad}{\mathrm{rad}}
\providecommand{\sgn}{}\renewcommand{\sgn}{\operatorname{sgn}}
\providecommand{\shA}{}\renewcommand{\shA}{\mathcal{A}}
\providecommand{\shC}{}\renewcommand{\shC}{\mathcal{C}}
\providecommand{\shE}{}\renewcommand{\shE}{\mathcal{E}}
\providecommand{\shJ}{}\renewcommand{\shJ}{\mathcal{J}}
\providecommand{\shK}{}\renewcommand{\shK}{\mathcal{K}}
\providecommand{\shL}{}\renewcommand{\shL}{L}
\providecommand{\shW}{}\renewcommand{\shW}{\mathcal{W}}
\providecommand{\sll}{}\renewcommand{\sll}{\mathfrak{sl}}
\providecommand{\so}{}\renewcommand{\so}{\mathfrak{so}}
\providecommand{\su}{}\renewcommand{\su}{\mathfrak{su}}
\providecommand{\supp}{}\renewcommand{\supp}{\operatorname{supp}}
\providecommand{\tr}{}\renewcommand{\tr}{\operatorname{tr}}
\providecommand{\uu}{}\renewcommand{\uu}{\mathfrak{u}}
\providecommand{\br}{}\renewcommand{\br}[1]{\langle #1\rangle}
\providecommand{\bra}{}\renewcommand{\bra}[1]{\langle #1|}
\providecommand{\file}{}\renewcommand{\file}[1]{\texttt{#1}}
\providecommand{\hp}{}\renewcommand{\hp}[2]{( #1 \,|\, #2 )}
\providecommand{\ip}{}\renewcommand{\ip}[2]{\langle #1, #2\rangle}
\providecommand{\ket}{}\renewcommand{\ket}[1]{|#1\rangle}
\providecommand{\rel}{}\renewcommand{\rel}[2]{D\bigl(#1 \,\|\, #2\bigr)}
\providecommand{\rev}{}\renewcommand{\rev}[1]{\widetilde{#1}}
\providecommand{\slashed}{}\renewcommand{\slashed}[1]{{#1}\mkern-10mu/}
\title{\bfseries One relation, read twice}
\author{Leon Butler\\ \normalsize\href{mailto:lbutler.math@gmail.com}{\texttt{lbutler.math@gmail.com}}\\ \normalsize ORCID: \href{https://orcid.org/0009-0006-4449-5617}{0009-0006-4449-5617}}
\date{Generated by the Transport Lab, 2 November 2026}
\begin{document}
\maketitle
\begin{abstract}
\normalsize
A Clifford algebra rests on one relation, which says how two vectors multiply when their order is swapped.

Written on an orthogonal basis it is a metric, the language of geometry; written on a Witt basis of paired null vectors it is the rule obeyed by fermions, the language of quantum theory. The geometric reading carries exactly one conversion constant, the speed of light.

Applied to the electron, the algebra ties the imaginary unit to the choice of rest frame. Granting stated premises, including that recording is irreversible and Hypothesis (S), measurement leaves every outcome in its own branch and nothing physically collapses. That stance predicts that no collapse signature will be found.
\end{abstract}
\newpage
\noindent \textbf{Authorship and disclosure. }The founding idea of the series is the author's: that division by zero in a Clifford algebra should be read not as a prohibition but as a transport into a higher-dimensional algebra, and that the transported division need not be accompanied by an inverse operation. The direction of investigation and the critical review of the results are likewise the author's. The mathematical development and the drafting of the text were carried out in extended collaboration with Claude Fable 5 (Anthropic), an AI system.

\noindent\fbox{\parbox{\dimexpr\linewidth-2\fboxsep-2\fboxrule}{\textbf{This is not a deposited paper.} It was generated by the Transport Lab on 2 November 2026 and assembled from approved parts of its store. It follows results across several papers of the series; each is stated with its proof and credited to the paper that first proved it. The official versions of those papers are on Zenodo, in the Transport Papers community, \url{https://zenodo.org/communities/transport-papers}.}}
\newpage

\section*{Introduction}
The walk starts with the two readings of the relation, follows the geometric reading to the speed of light, then the quantum reading through the electron's law to Schrödinger's equation, and ends with measurement and the prediction it makes.

\section{What we assume}

\begin{postulate}[CPT covariance of the kinematic structures and of every action]\label{c:dark-matter.premise-cpt-covariance}
This is an assumption, partly checked: a particular symmetry that swaps matter with antimatter, reverses time and mirrors space is taken to leave every basic structure of the framework, and every one of its laws including the strengths of the forces, unchanged. Worked out exactly, the symmetry involves complex conjugation, the same kind of operation quantum field theory uses for this symmetry, rather than an ordinary linear swap. The series verifies the assumption for the geometric structures and for gravity, and assumes it for the imported description of the forces and of matter. It is the first premise behind the claim that the matter-antimatter orientation of the universe is a free choice rather than something the framework fixes.
\end{postulate}

\begin{postulate}[The flag as datum: the sign of the Witt flag is configuration data]\label{c:dark-matter.premise-flag-as-datum}
The framework's origin event makes one yes-or-no choice, which fixes both which of matter and antimatter is which and the direction of time. This assumption says where that choice lives: it is a fact about our particular universe, set when the universe began, not a built-in feature of the theory's laws or structure. A different universe obeying the same laws could have made the opposite choice. Several results, including the claim that the choice is a starting condition rather than a setting of the laws, rest on this assumption. It would be shown wrong if some part of the theory's fixed structure were found to prefer one answer.
\end{postulate}

\begin{postulate}[Recording is irreversible, stated so it can be refuted]\label{c:one-relation.recording-irreversible-stated-so-it-can}
When an instrument records a measurement result, that record is assumed to be permanent: it cannot be undone. The series proves something similar for a more abstract process but not for real laboratory records, so this is stated openly as an assumption that could be refuted.
\end{postulate}

\section{One relation, read twice}

These results take the defining relation of a Clifford algebra and ask what follows from reading it as geometry and as quantum theory.

The defining relation of a Clifford algebra is \[ \{e\_{a}, e\_{b}\} = 2g\_{ab}\mathbf{1} \] which on an orthogonal basis displays a metric $g$, measuring lengths and angles.

On a Witt basis of paired null vectors the same relation reads $\{c(w_{i}), c(\bar{w}_{j})\} = \delta_{ij}$, exactly the canonical anticommutation relations of fermionic modes. So geometry and quantum structure are two charts on one algebra, at the level of structure rather than dynamics.

\begin{proposition}[The two presentations]\label{c:one-relation.two-presentations}
The basic rule of a Clifford algebra can be written in two ways. In one it describes distances and angles, the language of geometry and gravity; in the other it is exactly the rule obeyed by fermions, the language of quantum theory. The paper takes this to mean that geometry and quantum structure are two views of one object, though only at the level of structure, not of how things evolve.
\emph{First proved in} T16, Prop 3.1 (\href{https://doi.org/10.5281/zenodo.22723212}{doi:10.5281/zenodo.22723212}).
\end{proposition}
\noindent\emph{Rests on:} \ref{c:one-relation.imported-pure-spinors-incidence-chart}, \ref{c:std-algebra.vector-squares}, \ref{c:std-algebra.witt-basis-car}.

If geometry can be read as a system of fermion modes, one can count the modes.

There are two counts, a real capacity using real null vectors and a complex one allowing complex vectors, \[ r\_{\R}(p,q) = \min(p,q), \qquad r\_{\C}(p,q) = \lfloor d/2 \rfloor \] with $d = p + q$. For spacetime $(1,3)$ they are $1$ and $2$, and for the six-dimensional $(2,4)$ they are $2$ and $3$.

\begin{proposition}[The two capacities, and three consequences]\label{c:one-relation.two-capacities-three-consequences}
Reading spacetime geometry as a system of fermion modes, there are two ways to count the modes: using only real directions, or allowing complex ones. The two counts usually differ, and for ordinary spacetime they give one and two. Each step up the tower of extended spaces adds exactly one mode to both counts, and the complex count of the extended six-dimensional space is three.
\emph{First proved in} T16, Prop 3.2 (\href{https://doi.org/10.5281/zenodo.22723212}{doi:10.5281/zenodo.22723212}).
\end{proposition}
\noindent\emph{Rests on:} \ref{c:one-relation.imported-witts-extension-theorem}, \ref{c:one-relation.two-presentations}, \ref{c:tower.witt-index-quantum-capacity}, \ref{c:std-algebra.vector-squares}, \ref{c:std-algebra.witt-decomposition}.

The geometric reading carries something the quantum reading does not see.

Time and space are measured in independent units, so a quadratic form on events needs a conversion constant $c$, \[ Q = c\^{}{2}t\^{}{2} - |\mathbf{x}|\^{}{2} \] in the Lorentzian case, and the null cone is the invariant ratio it fixes. Such a finite ratio exists only when the signature is mixed, $pq \neq 0$. This places the constant without deriving its value.

\begin{lemma}[The third constant is the conversion constant between the charts]\label{c:one-relation.third-constant-conversion-constant-between-charts}
The speed of light can be seen as a conversion factor between units of time and units of length. This result shows such a factor exists only when spacetime mixes time-like and space-like directions, that there is only one such number, and that it is invisible in the quantum description and appears only on passing to the geometric one. It also places the three fundamental constants at three different levels of structure, without deriving any of their values.
\emph{First proved in} T16, Lem 4.4 (\href{https://doi.org/10.5281/zenodo.22723212}{doi:10.5281/zenodo.22723212}).
\end{lemma}
\noindent\emph{Rests on:} \ref{c:one-relation.imported-null-cone-determines-form-zeemans}, \ref{c:one-relation.imported-witts-extension-theorem}, \ref{c:std-algebra.bivector-exponential}.

The quantum reading now turns to the electron, whose relativistic law is first order.

A law $\nabla\psi = m\,\psi\,C$ with a constant factor $C$, applied twice, gives the Klein-Gordon equation \[ \nabla\^{}{2}\psi = -m\^{}{2}\psi \] exactly when $C^{2} = -1$. With $C^{2} = +1$ the sign is tachyonic, and a non-scalar square gives no scalar equation at all.

\begin{lemma}[The form of the law is forced by its square]\label{c:one-relation.form-law-forced-square}
A first-order wave equation of this type, applied twice, should give the familiar relativistic relation between energy, momentum and mass. Whether it does depends only on the square of one constant factor in the law. That square must be minus one; any other choice gives faster-than-light behaviour or no sensible equation at all.
\emph{First proved in} T16, Lem 5.4 (\href{https://doi.org/10.5281/zenodo.22723212}{doi:10.5281/zenodo.22723212}).
\end{lemma}

That condition can be solved completely.

The odd square roots of minus one in $\Cl(1,3)$ are \[ C = (\cos t + \sin t\, I)\,u \] with $u$ a unit spacelike vector and $I$ the pseudoscalar. A constant rotor moves $u$ and the chiral phase removes $t$, so for a single free fermion with minimal electromagnetic coupling every allowed law is the Dirac-Hestenes equation up to a constant field redefinition.

\begin{lemma}[Odd square roots of $-1$, and the chiral phase]\label{c:one-relation.odd-square-roots-chiral-phase}
An earlier result says the constant factor in the electron's law must square to minus one. This one lists every allowed factor and shows that, for a single free electron, they all give the same equation seen in a rotated frame or with a harmless constant phase. Once other interactions are included, that phase is no longer harmless.
\emph{First proved in} T16, Lem 5.5 (\href{https://doi.org/10.5281/zenodo.22723212}{doi:10.5281/zenodo.22723212}).
\end{lemma}
\noindent\emph{Rests on:} \ref{c:std-algebra.pseudoscalar-commutation}, \ref{c:primer.sign}, \ref{c:std-algebra.vector-squares}.

Quantum mechanics needs an imaginary unit and slow-motion physics needs a rest frame, and these look like two separate choices.

With the law aligned to $C = Iu$, the admissible imaginary units $J$ correspond one to one with the unit timelike vectors orthogonal to $u$, through \[ \gamma\_{0} = CJ \] so picking one picks the other. The imaginary unit is not an extra ingredient put in by hand.

\begin{lemma}[One datum, not two: the imaginary unit is the rest frame read through the law]\label{c:one-relation.one-datum-not-two-imaginary-unit}
Quantum mechanics uses an imaginary unit, and nonrelativistic physics uses a rest frame. In this algebraic form of the electron's equation, once the law is fixed, picking one automatically picks the other. So the imaginary unit is not an extra ingredient that has to be put in by hand.
\emph{First proved in} T16, Lem 5.6 (\href{https://doi.org/10.5281/zenodo.22723212}{doi:10.5281/zenodo.22723212}).
\end{lemma}
\noindent\emph{Rests on:} \ref{c:one-relation.odd-square-roots-chiral-phase}, \ref{c:std-algebra.duality}.

With the law and its imaginary unit in place, the slow-motion limit can be taken.

The operator is the vector derivative $\nabla = \gamma^{\mu}\partial_{\mu}$, forced, and the law is the Dirac-Hestenes equation \[ \nabla\psi\, I\sigma\_{3} = m\,\psi\gamma\_{0} \] whose slow limit is Schrödinger's equation, at a structural level and for a single free fermion.

The orientation of time comes from an earlier argument that assumes Hypothesis (S), and the rest frame itself is supplied by the laboratory, not derived.

\begin{proposition}[Schrödinger as the slow shadow]\label{c:one-relation.schr-odinger-slow-shadow}
The everyday equation of quantum mechanics, Schr\"odinger's, is shown to be the slow-motion limit of the relativistic electron equation, which is itself built from the basic algebra of spacetime. The paper tracks where each ingredient comes from, at a structural level: the mass from an earlier mechanism, at the observation level that mechanism's reading holds; the orientation half of the direction of time from an earlier argument that holds if the slots behave like the pairs (Hypothesis (S)); and the imaginary unit from the choice of rest frame. Only the first-order form of the law and the standard slow-motion expansion are imported, for a single free fermion. The rest frame itself is supplied by the laboratory, not derived.
\emph{First proved in} T16, Prop 5.2 (\href{https://doi.org/10.5281/zenodo.22723212}{doi:10.5281/zenodo.22723212}).
\end{proposition}
\noindent\emph{Rests on:} \ref{c:one-relation.arrow-frequency-are-one-bit-level}, \ref{c:one-relation.form-law-forced-square}, \ref{c:one-relation.imported-flag-arrow-superselection}, \ref{c:one-relation.imported-condensate-derived-content}, \ref{c:one-relation.odd-square-roots-chiral-phase}, \ref{c:one-relation.one-datum-not-two-imaginary-unit}, \ref{c:one-relation.vector-derivative-covariant-first-order-operator}, \ref{c:standard-model.hyp-s-slot-statistics}, \ref{c:std-algebra.vector-squares}.

Quantum mechanics also needs a rule for chances, Born's rule.

The unique conserved density is the time component of the current \[ J = \psi\gamma\_{0}\tilde{\psi} \] and under the indifference premise, that nothing may prefer one of several equivalent outcomes, it is the only consistent assignment of weights. The result is conditional on that premise, and the density is relative to the laboratory's rest frame.

\begin{proposition}[The measure from indifference: Born's rule located]\label{c:one-relation.measure-indifference-borns-rule-located}
Born's rule says the chance of a quantum outcome is the squared size of its part of the wave. Here the relevant density is fixed as the unique conserved, positive one, measured in the laboratory's rest frame. Then the series' indifference premise, a declared starting point and not a theorem, that nothing may prefer one of several equivalent outcomes, forces equal chances for equal parts and builds up to the squared-size rule. One assumption remains, that chances exist and add up, and the paper names it as what is still postulated, so the result is conditional on both.
\emph{First proved in} T16, Prop 6.2 (\href{https://doi.org/10.5281/zenodo.22723212}{doi:10.5281/zenodo.22723212}).
\end{proposition}
\noindent\emph{Rests on:} \ref{c:one-relation.dirac-current-conserved-density}, \ref{c:one-relation.imported-indifference-instrument}.

The same premises bear on what happens in a measurement.

Granting indifference, that branches cannot see one another, and that recording is irreversible, with Hypothesis (S) behind the count, every outcome happens in its own branch and nothing physically collapses. Picking one outcome would be an uncaused preference.

Why this branch is the one experienced is left open.

\begin{proposition}[The record chain votes: measurement without collapse]\label{c:one-relation.record-chain-votes-measurement-without-collapse}
When a measurement is made, the framework, granting its stated premises (that an unexplained preference is forbidden, that recording is irreversible, and, behind its count of free numbers, Hypothesis (S) on the slots' statistics), says every possible outcome happens, each in its own branch that cannot see the others, rather than one outcome being picked by a physical collapse. It argues this in two ways: picking one outcome would be an unexplained preference, and collapse would need an endless supply of new random choices the framework does not have. The parallel with the cosmic case holds only in the limit of ideal recording, since laboratory decoherence is approximate, and the account leaves open why this branch is the one experienced. This is a testable, conditional bet: any confirmed sign of real wavefunction collapse would refute it.
\emph{First proved in} T16, Prop 6.5 (\href{https://doi.org/10.5281/zenodo.22723212}{doi:10.5281/zenodo.22723212}).
\end{proposition}
\noindent\emph{Rests on:} \ref{c:one-relation.imported-flag-arrow-superselection}, \ref{c:one-relation.imported-indifference-instrument}, \ref{c:one-relation.measure-indifference-borns-rule-located}, \ref{c:standard-model.trichotomy-what-kind-number}, \ref{c:dark-matter.asymmetry-arrow-counted-fifteenth-refutation}, \ref{c:dark-matter.orientation-arrow-seventeenth-refutation}, \ref{c:standard-model.hyp-s-slot-statistics}, \ref{c:one-relation.recording-irreversible-stated-so-it-can}, \ref{c:dark-matter.premise-cpt-covariance}, \ref{c:dark-matter.premise-flag-as-datum}.

A stance on measurement is worth most if it can fail.

The framework predicts no objective-collapse signature: no spontaneous-localisation noise and no gravitationally induced collapse. A single confirmed detection would refute the prediction together with the argument behind it.

\begin{prediction}[No objective-collapse signature]\label{c:pred.no-objective-collapse}
Some proposed theories say quantum superpositions collapse on their own, perhaps because of gravity, and these would leave tiny detectable traces such as extra random noise. This framework predicts that no such traces exist at any level of sensitivity. The prediction rests on the framework's account of measurement, under that account's premises. Experiments so far have found none, and a single confirmed detection would refute the prediction together with the argument behind it.
\emph{Kill condition:} One confirmed collapse signature.
\end{prediction}

\section{Results used from elsewhere}
The proofs above use the following results of the series, stated briefly; each is proved in the paper named.

\begin{fact}[The asymmetry: the arrow counted, and the fifteenth refutation]\label{c:dark-matter.asymmetry-arrow-counted-fifteenth-refutation}
The framework inherits rather than generates its asymmetry, which is confined to this tower and the mirror, with a common value under premise P4; a constant early $\theta$-angle as its source is excluded by the neutron dipole bound, and, on a flagged stance, the value $\epsilon$ is assumed to be initial data of the descent.
\emph{Source:} T14, Prop 5.1 (\href{https://doi.org/10.5281/zenodo.22722953}{doi:10.5281/zenodo.22722953}).
\end{fact}

\begin{fact}[The orientation is the arrow: the seventeenth refutation]\label{c:dark-matter.orientation-arrow-seventeenth-refutation}
Under CPT covariance (an antilinear conjugation) of the kinematic structures and of every action (gauge couplings included), Hypothesis (S), and the premise that the flag is a datum of each solution, the particle-antiparticle orientation takes both values on every solution space, so it is initial data (not a parameter of the laws), one choice with the arrow of time.
\emph{Source:} T14, Prop 5.2 (\href{https://doi.org/10.5281/zenodo.22722953}{doi:10.5281/zenodo.22722953}).
\end{fact}

\begin{fact}[The arrow and the frequency are one bit, at $\Cl(1,3)$ level]\label{c:one-relation.arrow-frequency-are-one-bit-level}
Given the fixed Dirac-Hestenes law, the frequency sign and the direction of the rest frame's time are a single bit, since $C(-J) = -(CJ)$.
\emph{Source:} T16, Cor 5.7 (\href{https://doi.org/10.5281/zenodo.22723212}{doi:10.5281/zenodo.22723212}).
\end{fact}

\begin{fact}[The Dirac current is the conserved density]\label{c:one-relation.dirac-current-conserved-density}
Among constant sesquilinear currents, the Dirac current $\bar{\psi}\gamma^{\mu}\psi$ is the unique conserved one up to scale, with positive density $\psi^{\dagger}\psi$.
\emph{Source:} T16, Lem 6.1 (\href{https://doi.org/10.5281/zenodo.22723212}{doi:10.5281/zenodo.22723212}).
\end{fact}

\begin{fact}[Imported: the condensate and the derived content]\label{c:one-relation.imported-condensate-derived-content}
The condensate's bracket is fixed by the join, condensation is an observation-level reading under an independence assumption, the democratic alignment rests on the indifference premise, and on a vacuum Einstein background mass is the grip on the frozen point with $m^{2} = -w(w+3)\Lambda/3$ at Klein-Gordon level, the gauge content coming from the forced connection.
\emph{Source:} T16, Fact 2.6 (\href{https://doi.org/10.5281/zenodo.22723212}{doi:10.5281/zenodo.22723212}).
\end{fact}

\begin{fact}[Imported: the flag, the arrow, and superselection]\label{c:one-relation.imported-flag-arrow-superselection}
Imported: the descent's single Witt-flag orientation bit is both the arrow of time and the particle-antiparticle orientation, is initial data under three stated premises, CPT covariance, Hypothesis (S) and the flag as a datum, and tower charge is exactly superselected.
\emph{Source:} T16, Fact 2.5 (\href{https://doi.org/10.5281/zenodo.22723212}{doi:10.5281/zenodo.22723212}).
\end{fact}

\begin{fact}[Imported: the indifference instrument]\label{c:one-relation.imported-indifference-instrument}
Indifference, a structural premise and not a theorem, equates weights along exact symmetry orbits, leaves genuine moduli as data (an observation-level reading), and treats which of a mirror pair is ours as indexical (P2).
\emph{Source:} T16, Fact 2.7 (\href{https://doi.org/10.5281/zenodo.22723212}{doi:10.5281/zenodo.22723212}).
\end{fact}

\begin{fact}[Imported: the null cone determines the form, and Zeeman's theorem]\label{c:one-relation.imported-null-cone-determines-form-zeemans}
Imported: a mixed-signature null cone fixes its quadratic form up to scale, and for $n \ge 3$ the causal automorphisms of Minkowski space are the conformal ones (Zeeman).
\emph{Source:} T16, Fact 2.9 (\href{https://doi.org/10.5281/zenodo.22723212}{doi:10.5281/zenodo.22723212}).
\end{fact}

\begin{fact}[Imported: pure spinors and the incidence chart]\label{c:one-relation.imported-pure-spinors-incidence-chart}
Imported: every conformal spinor is a pure spinor, the Fock vacuum of its own three-pair system, and spacetime points are null vectors of $\C^6$, equivalently lines of the projective spinor space.
\emph{Source:} T16, Fact 2.3 (\href{https://doi.org/10.5281/zenodo.22723212}{doi:10.5281/zenodo.22723212}).
\end{fact}

\begin{fact}[Imported: Witt's extension theorem]\label{c:one-relation.imported-witts-extension-theorem}
Over a field of characteristic not two, every isometry between subspaces of a nondegenerate quadratic space extends to the whole space, so the isometry group is transitive on subspaces of each isometry type.
\emph{Source:} E. Witt, Theorie der quadratischen Formen in beliebigen K"orpern, J. Reine Angew. Math. 176 (1937) 31--44.; T. Y. Lam, Introduction to Quadratic Forms over Fields, Graduate Studies in Mathematics 67, American Mathematical Society (2005)., Ch.~I, \S4.
\end{fact}

\begin{fact}[The vector derivative is the covariant first-order operator]\label{c:one-relation.vector-derivative-covariant-first-order-operator}
Every spin-covariant, constant-coefficient, first-order operator acting by left multiplication on $\Cl(1,3)$-valued fields is an invertible constant times the vector derivative $\nabla = \gamma^{\mu}\partial_{\mu}$.
\emph{Source:} T16, Lem 5.1 (\href{https://doi.org/10.5281/zenodo.22723212}{doi:10.5281/zenodo.22723212}).
\end{fact}

\begin{fact}[The sign of $\omega^2$]\label{c:primer.sign}
The pseudoscalar of $\Cl(p,q)$ satisfies $\omega^{2} = (-1)^{(p-q)(p-q-1)/2}$, so $\omega^{2} = +1$ exactly when $p - q \equiv 0$ or $1 \pmod 4$.
\emph{Source:} T0, Lem 9.5 (\href{https://doi.org/10.5281/zenodo.22714501}{doi:10.5281/zenodo.22714501}).
\end{fact}

\begin{fact}[Hypothesis (S): slot statistics, derived from the slot-module hypothesis]\label{c:standard-model.hyp-s-slot-statistics}
Under the slot-module hypothesis, the composite's slots are fermionic modes obeying exclusion and forming one Clifford system with its pairs, its state space is the five-mode Fock space $\Lambda^{\bullet}(\C^{2} \oplus W)$ with $U(2) \times U(3)$ acting as $\Lambda^{\bullet}(A \oplus B)$, chirality is the parity of the mode number, and slots of one tower parity whose grading is the chirality are odd.
\emph{Source:} T12, Hypothesis (S) (\href{https://doi.org/10.5281/zenodo.22722669}{doi:10.5281/zenodo.22722669}).
\end{fact}

\begin{fact}[The trichotomy: what kind of number $N$ is]\label{c:standard-model.trichotomy-what-kind-number}
The number $N = M_{P}/H$ rescales while every derived discrete invariant stays fixed (proved), which is read at observation level as $N$ not being kinematic; it is imported as not dynamical, and under the action as written it is a parameter of the laws, becoming initial data only if the norm is promoted to a field.
\emph{Source:} T12, Prop 7.6 (\href{https://doi.org/10.5281/zenodo.22722669}{doi:10.5281/zenodo.22722669}).
\end{fact}

\begin{fact}[Exponentials of bivectors with scalar square]\label{c:std-algebra.bivector-exponential}
If $A^{2} = \lambda \in \R$ then $e^{A}$ is $\cos\theta + \frac{\sin\theta}{\theta}A$, $\cosh\theta + \frac{\sinh\theta}{\theta}A$ or $1 + A$, a rotor with $e^{A}\widetilde{e^{A}} = 1$.
\emph{Source:} the algebra core of the Transport Lab, proved there with an exact check; standard, see C. Doran and A. Lasenby, Geometric Algebra for Physicists, Cambridge University Press (2003), Ch. 4 (foundations of geometric algebra: rotors, versors).
\end{fact}

\begin{fact}[Duality]\label{c:std-algebra.duality}
$X \mapsto XI$ sends $\Cl^{k} \to \Cl^{n-k}$ with $a \rfloor (XI) = (a \wedge X)I$, $a \wedge (XI) = (a \rfloor X)I$; for $r = 0$ it is an isomorphism and $[AI] = [A]^{\perp}$.
\emph{Source:} the algebra core of the Transport Lab, proved there with an exact check; standard, see C. Doran and A. Lasenby, Geometric Algebra for Physicists, Cambridge University Press (2003), Ch. 4 (foundations of geometric algebra).
\end{fact}

\begin{fact}[How the pseudoscalar commutes]\label{c:std-algebra.pseudoscalar-commutation}
$I X_k = (-1)^{k(n-1)} X_k I$; $I\tilde{I} = \prod_i Q(e_i)$, and $I^{2} = 0$ in a degenerate algebra.
\emph{Source:} the algebra core of the Transport Lab, proved there with an exact check; standard, see C. Doran and A. Lasenby, Geometric Algebra for Physicists, Cambridge University Press (2003), Ch. 4 (foundations of geometric algebra).
\end{fact}

\begin{fact}[The product of two vectors]\label{c:std-algebra.vector-squares}
$v^{2} = Q(v)$ and $uv + vu = 2B(u, v)$, so $uv = B(u, v) + \tfrac12(uv - vu)$.
\emph{Source:} the algebra core of the Transport Lab, proved there with an exact check; standard, see C. Doran and A. Lasenby, Geometric Algebra for Physicists, Cambridge University Press (2003), Ch. 4 (foundations of geometric algebra).
\end{fact}

\begin{fact}[Null pairs satisfy the canonical anticommutation relations]\label{c:std-algebra.witt-basis-car}
For a Witt basis ($B(v_i,v_j) = B(w_i,w_j) = 0$, $B(v_i,w_j) = \tfrac12\delta_{ij}$), $\{v_i,v_j\} = \{w_i,w_j\} = 0$ and $\{v_i,w_j\} = \delta_{ij}$ in $\Cl(V,Q)$.
\emph{Source:} the algebra core of the Transport Lab, proved there with an exact check; standard, see C. Doran and A. Lasenby, Geometric Algebra for Physicists, Cambridge University Press (2003), Ch. 9 (multiparticle states).
\end{fact}

\begin{fact}[Hyperbolic pairs and the Witt decomposition]\label{c:std-algebra.witt-decomposition}
A null vector in a nondegenerate plane has a unique hyperbolic partner, and $V = H_1 \perp \dots \perp H_{\min(p,q)} \perp A \perp \operatorname{rad} V$ with $A$ definite; maximal totally null subspaces have dimension $\min(p,q) + r$.
\emph{Source:} the algebra core of the Transport Lab, proved there with an exact check; standard, see T. Y. Lam, Introduction to Quadratic Forms over Fields, Graduate Studies in Mathematics 67, American Mathematical Society (2005), Ch. I (hyperbolic planes, the Witt decomposition, Witt's theorems).
\end{fact}

\begin{fact}[Witt index is quantum capacity]\label{c:tower.witt-index-quantum-capacity}
For a nondegenerate real quadratic space of signature $(p,q)$, the maximal number of fermionic mode pairs realisable by vectors of its Clifford algebra is the Witt index $\min(p,q)$, which rises by exactly one at each rung of the tower.
\emph{Source:} T3, Thm 4.3 (\href{https://doi.org/10.5281/zenodo.22721556}{doi:10.5281/zenodo.22721556}).
\end{fact}

\section*{Conclusion}
One relation, read twice, gives both the geometry of spacetime and the fermion structure of quantum theory.

The steps on measurement rest on the postulates set out under what we assume, that laboratory records are never unwritten and Hypothesis (S) on slot statistics, together with the indifference premise. A single confirmed collapse signature would refute the prediction together with the argument behind it.

\end{document}